\chapter{Weak balance in the complete circle packing}
\label{app:weak-balance}

The strict balance condition in Appendix~\ref{app:balanced-terminal-fans}
is sufficient to control the complementary part of the invariant. We now
allow equality in that condition. The change is useful because the face
weights used in a coefficient estimate need not remain the curvatures
of a circle packing. Moving them a small distance from the complete
packing creates a first-order local gain. Stationarity makes the cost
to the complete action only second order. We account for both changes
on the same shared face variables.

\section{The unchanged diagram and the volume statement}

Fix a finite simple triangulation of the sphere with at least four
vertices, a perfect matching of its cubic dual, and one nonmatching
edge as the reduced mark. Use the augmented diagram and its labeled
crossing disks constructed in Appendix~\ref{app:morse-geometric-contribution}.
There are $R$ matched pairs of triangles and $R+2$ white vertices.
The marked edge joins the protected vertices $u,v$; let $V$ be the
set of the $R$ ordinary vertices. The complete packing has curvatures
$w_u^*=w_v^*=0$ and $w_f^*>0$ for $f\in V$. Curvature is the
reciprocal radius of a finite circle. Write
\[
 q_f=\log w_f^*,\qquad
 \Phi=\frac{\Vol(Y_D)}{2\pi},
\]
where $Y_D$ is the actual augmented link complement.

At a crossing region the cyclic face order is $(x,b,y,c)$ and the
matched diagonal is $xy$. The polynomial arguments are ordered
$(x,b,c,y)$. To state the condition, define
\[
 \Lambda(s)=-\int_0^s\log|2\sin z|\,\dd z,
 \qquad f(a)=\frac{\Lambda(\pi a)+\Lambda(\pi/2-\pi a)}{\pi}
       \quad(0\le a\le\tfrac12).
\]
For a triple of nonnegative weights $w=(w_1,w_2,w_3)$ with at least
one positive entry, its triangle potential is
\[
 Q(w)=\max_{\substack{a_i\ge0\ (w_i>0)\\
                     \sum_{i:w_i>0}a_i=1/2}}
           \sum_{i:w_i>0}\{f(a_i)+a_i\log w_i\}.
\]
The zero-weight corners are omitted from this maximization. If only
$w_i$ is positive, then $Q(w)=\tfrac12\log w_i$. These conventions
avoid taking a logarithm of a protected zero. Put
\[
 X_e=Q(w_x^*,w_b^*,w_y^*)+Q(w_x^*,w_c^*,w_y^*).
\]
Only positive middle curvatures occur inside logarithms below.

\begin{theorem}[Volume limit under weak balance]\label{wb:volume}
Suppose the actual complete packing satisfies
\begin{equation}\label{wb:balance}
 X_e\ge\log w_b^*\quad(w_b^*>0),\qquad
 X_e\ge\log w_c^*\quad(w_c^*>0)
\end{equation}
at every crossing region. Fix retained half twists
$\epsilon_e\in\{0,\pm1\}$ whose original strand connections give
one knot, with no additional closed main component. There is a
threshold $T$, depending on these fixed data, such that independently
chosen fixed integers $p_e$ with $|p_e|>T$ give a hyperbolic knot
$K_{\boldsymbol p}$ and
\[
 \lim_{N\to\infty}\frac{2\pi}{N}
       \log|\langle K_{\boldsymbol p}\rangle_N|
       =\Vol(S^3\setminus K_{\boldsymbol p}).
\]
The crossing powers are $n_e=\epsilon_e-2p_e$. The invariant is the
original reduced, zero-framed Kashaev invariant with unknot value one,
and the limit includes both parities and composite levels. The
diagram, its mark, and all powers are fixed before $N$ varies.
\end{theorem}

No terminal tangles are inserted in this theorem. Its restriction is
the packing inequality~\eqref{wb:balance}; it is not a theorem for all
unchanged matched spheres. A quantified extension below permits some
positive defects in that inequality. The proof uses the full geometric
contribution of Appendix~\ref{app:morse-geometric-contribution} and the
finite norm bound of Appendix~\ref{app:uniform-filling-bound}.

\section{Triangle entropy, its Hessian, and the complete action}

For an original triangle $t$, let $I_t\subset V$ be its ordinary
corners. This set is nonempty. Put
\begin{equation}\label{wb:triangle}
\begin{split}
 f(a)&=\frac{\Lambda(\pi a)+\Lambda(\pi/2-\pi a)}{\pi},
       \qquad 0\le a\le\tfrac12,\\
 Q_t(q')&=\max_{\substack{a_i\ge0\ (i\in I_t)\\
                         \sum_{i\in I_t}a_i=1/2}}
               \sum_{i\in I_t}\bigl(f(a_i)+a_iq'_i\bigr),
\end{split}
\end{equation}
where $\Lambda(s)=-\int_0^s\log|2\sin z|\,\dd z$.
Protected corners are deleted from the supported simplex; their weights
remain exactly zero. The notation $Q(w_x,w_b,w_y)$ means this function
with $q'_i=\log w_i$ at the positive coordinates. This is the absolute
normalization of the cofactor action, including the determinant, beta
and half-index factors. In particular a singleton row has
$Q_t(q')=q'_i/2$.

Direct differentiation gives
\[
 f(0)=f(1/2)=0,\qquad f'(a)=\log\cot(\pi a),\qquad
 f''(a)=-\frac{2\pi}{\sin(2\pi a)}<0.
\]
Thus $f(a)>0$ for $0<a<1/2$. If two or three weights are positive,
write $w_i=e^{q'_i}$ and
\begin{equation}\label{wb:angles}
 D_t^2=\sum_{\{i,j\}\subset I_t}w_iw_j,
 \qquad a_{ti}=\frac1\pi\arctan\frac{w_i}{D_t}.
\end{equation}
These loads sum to $1/2$. For three weights this follows from
$D_t^2-w_iw_j=w_k(w_i+w_j)$ and the addition formula for tangent;
for two weights the two tangent ratios are reciprocal. Moreover,
$f'(a_{ti})+q'_i=\log D_t$. Strict concavity proves that these are
the unique maximizing loads in~\eqref{wb:triangle}. Consequently
\[
 \partial_{q'_i}Q_t=a_{ti},\qquad
 E_t(q'):=Q_t(q')-q'\cdot\nabla Q_t(q')
        =\sum_{i\in I_t}f(a_{ti}).
\]
The row entropy $E_t$ is positive on a nonsingleton row and zero on
a singleton row.

For distinct ordinary corners, differentiating~\eqref{wb:angles} gives
\begin{equation}\label{wb:hessian}
 \partial_{q'_iq'_j}Q_t=-\kappa_{ij}^t,
 \qquad
 \kappa_{ij}^t=\frac{w_iw_j}{2\pi D_t(w_i+w_j)}
                    \le\frac1{4\pi}.
\end{equation}
For three corners, the cancellation uses
$D_t^2+w_i^2=(w_i+w_j)(w_i+w_k)$. The upper bound follows from
$D_t\ge\sqrt{w_iw_j}$ and $w_i+w_j\ge2\sqrt{w_iw_j}$.
Homogeneity gives zero row sums in the Hessian. Its quadratic form is
therefore
\[
 h^{\mathsf t}\nabla^2Q_t(q')h
       =\sum_{\{i,j\}\subset I_t}\kappa_{ij}^t(h_i-h_j)^2.
\]
This also proves positive semidefiniteness directly.

Set
\begin{equation}\label{wb:action}
 P(q')=\sum_tQ_t(q')-\sum_{f\in V}q'_f.
\end{equation}
At the complete packing $q$, the face-star equations and calibrated
geometric action of Appendix~\ref{app:morse-geometric-contribution} give
\begin{equation}\label{wb:stationary}
 \sum_{t\ni f}a_{tf}=1,\qquad \nabla P(q)=0,
 \qquad P(q)=\Phi.
\end{equation}
These are statements about the complete geometric contribution. They
do not assume a bound for the complementary part of the invariant.

\section{Changing the common weights}

For $0<t<1$, replace every ordinary weight by
\[
             w_f(t)=(w_f^*)^{1-t}.
\]
Keep the two protected weights zero. These numbers specify one valid
Cauchy coefficient torus. They need not describe another packing.
Define
\begin{equation}\label{wb:energy}
 \mathcal E=\frac1{4\pi}
       \sum_t\sum_{\{i,j\}\subset I_t}(q_i-q_j)^2
   =\frac1{2\pi}\sum_{\substack{ij\text{ a white edge}\\i,j\in V}}
                         (q_i-q_j)^2.
\end{equation}
The second expression counts the two triangles incident to each
ordinary white edge. Taylor's formula, \eqref{wb:hessian}, and
\eqref{wb:stationary} give the entire action cost
\begin{equation}\label{wb:cost}
 0\le P((1-t)q)-\Phi
   =\int_0^t(t-s)q^{\mathsf t}\nabla^2P((1-s)q)q\,\dd s
   \le\tfrac12t^2\mathcal E.
\end{equation}

The local gain is first order. Put
$E_t(\beta)=Q_t(\beta q)-\beta q\cdot\nabla Q_t(\beta q)$ for
$0<\beta\le1$. Differentiation gives
\[
 E_t'(\beta)=-\beta q^{\mathsf t}\nabla^2Q_t(\beta q)q\le0.
\]
Integrating the derivative of $Q_t(\beta q)/\beta$ yields
\begin{equation}\label{wb:entropy-gain}
 Q_t(\beta q)=\beta Q_t(q)+
     \beta\int_\beta^1\frac{E_t(s)}{s^2}\,\dd s
 \ge\beta Q_t(q)+(1-\beta)E_t(q).
\end{equation}
Singleton rows satisfy the same identity, with zero entropy.

At a crossing region, abbreviate $Q_b=Q(x,b,y)$ and $Q_c=Q(x,c,y)$.
The three local rates, uniform in face phases and throughout the
real spectral interval, are
\begin{equation}\label{wb:rates}
 A_C=X=Q_b+Q_c,\quad
 A_L=\max(\log c,Q_b+\tfrac12\log c),\quad
 A_R=\max(\log b,Q_c+\tfrac12\log b).
\end{equation}
A left or right label with protected overloaded middle is omitted:
its finite polynomial is identically zero. The source estimates
behind~\eqref{wb:rates} are the coefficient bound for the central
cofactor and the complete unused-prefix and unused-suffix estimate
for a one-sided path, as used in
Section~\ref{bl:every-sector-containing-a-one-sided-crossing-circle-has-a-global-spectral-gap}.
They hold at arbitrary fixed positive ordinary weights. Their derivation
precedes any use of complete balance.

For a surviving side $(e,s)$, let $z$ be its overloaded middle and let
$t(e,s)$ be the triangle containing $z$. Define, at the complete packing,
\begin{equation}\label{wb:defects}
 \epsilon_{e,s}=\max(0,q_z-X_e(q)),\qquad
 h_{e,s}=E_{t(e,s)}(q)>0.
\end{equation}
The strict positivity needs a geometric incidence observation. The
outer edge is matched, whereas the protected edge is nonmatching.
Since the sphere is simple, the outer pair cannot be precisely the
two protected vertices. Hence the triangle containing a positive
overloaded middle has another ordinary corner.

For the left side the exact loss is
\[
 X-A_L=\min(X-\log c,Q_c-\tfrac12\log c).
\]
At completeness the second expression is positive: the pure-middle
simplex vertex has value $\tfrac12\log c$, while the unique maximizer
has positive load at every allowed corner. Apply
\eqref{wb:entropy-gain} to both expressions and put $\beta=1-t$.
The first receives the entropy of both triangles, while the second
receives that of the overloaded triangle. This proves, for either side,
\begin{equation}\label{wb:loss}
 (X_e-A_s)((1-t)q)
       \ge t h_{e,s}-(1-t)\epsilon_{e,s}.
\end{equation}
All weights in this comparison belong to the same global vector.

\section{A quantitative criterion and the full complementary sum}

Suppose for some $0<t<1$ that
\begin{equation}\label{wb:criterion}
 \gamma_0:=\min_{(e,s)\text{ surviving}}
        \{t h_{e,s}-(1-t)\epsilon_{e,s}\}
                    -\tfrac12t^2\mathcal E>0.
\end{equation}
This is a sufficient condition even when weak balance fails. If no
one-sided label survives, the complement is exactly zero and the
criterion is unnecessary. Otherwise put $h_* =\min h_{e,s}>0$.
Under weak balance all defects vanish. Choose
\begin{equation}\label{wb:choice}
 t_0=\begin{cases}
       \min(1/2,h_*/\mathcal E),&\mathcal E>0,\\
       1/2,&\mathcal E=0.
     \end{cases}
\end{equation}
Then $\gamma_0\ge t_0h_*/2>0$; when $\mathcal E=0$ the stronger
value is $h_*/2$. This proves that any number of equalities in
\eqref{wb:balance} cause no exceptional sector. More generally,
$\max\epsilon_{e,s}<t_0h_*/2$ suffices, with gap at least
$t_0h_*/2-\max\epsilon_{e,s}$.

Here is the passage from these real functions to the original invariant.
Put $d=N-1$ and write $S_N(\boldsymbol\sigma)$ for the exact shared
spectral polynomial of the Morse source. Its local numerator is
\[
 V_N(\sigma;\Omega)=
 \frac{\chi_N(\sigma)-\chi_N(\Omega)}{\sigma-\Omega},
 \qquad \chi_N(z)=2T_N(z/2)+2,
\]
where division is polynomial division and
$T_N(\cos\theta)=\cos(N\theta)$. Thus each spectral degree is at
most $d$, including the derivative data at repeated roots.
The shared polynomial is
\[
 S_N=[\prod_{f\in V}L_f^d]\,
       \left.\prod_eH_{V_N}
        (L_{x_e},L_{b_e},L_{c_e},L_{y_e};\sigma_e)
       \right|_{L_u=L_v=0}.
\]
The original marked cut, face measure and even-level coset are included
in this identity. No further factor $N^R$ is inserted.

Let $C_N$ select the central path sector at every crossing region,
with the odd central row included once, and let $D_N=S_N-C_N$.
A genuine label string in $D_N$ contains at least one left or right
choice. Denote its one-sided positions by $J$. Cauchy's coefficient
inequality on the single torus $|L_f|=(w_f^*)^{1-t}$, together with
\eqref{wb:cost} and~\eqref{wb:loss}, bounds its logarithmic rate by
\[
 \Phi+\tfrac12t^2\mathcal E
    -\sum_{(e,s)\in J}
                  \{t h_{e,s}-(1-t)\epsilon_{e,s}\}
 \le\Phi-\gamma_0.
\]
Condition~\eqref{wb:criterion} makes every summand in braces
positive, so other one-sided choices cannot undo the loss. The local
bounds are uniform in all face phases and real spectral values.
The original target is $d$, and $d/N=1-1/N$ changes the normalized
logarithm by only $O(1/N)$ at these fixed weights. There are at most
$3^R-1$ strings, a number independent of $N$. Hence
\begin{equation}\label{wb:remainder}
 \limsup_{N\to\infty}\frac1N
         \log\|D_N\|_{[-2,2]^R}\le\Phi-\gamma_0.
\end{equation}
The statement includes both level parities and retains all shared
face targets.

\begin{figure}[ht]
\centering
\begin{tikzpicture}[x=6.8cm,y=4.6cm,>=Latex,font=\small]
  \definecolor{gainblue}{HTML}{1B6780}
  \definecolor{costamber}{HTML}{AD601A}
  \fill[gainblue!10] plot[domain=0:.82,samples=60] (\x,\x)
    -- plot[domain=.82:0,samples=60] (\x,{.5*\x*\x}) -- cycle;
  \draw[->] (0,0)--(1.06,0) node[right] {$t$};
  \draw[->] (0,0)--(0,.96) node[above,align=center] {change in\\logarithmic rate};
  \draw[gainblue,very thick] plot[domain=0:.84] (\x,\x);
  \draw[costamber,very thick] plot[domain=0:1,samples=60] (\x,{.5*\x*\x});
  \node[gainblue,anchor=west] at (.84,.84) {$t h_*$};
  \node[costamber,anchor=west] at (1,.5) {$\tfrac12 t^2\mathcal E$};
  \draw[dashed,gray] (.5,0)--(.5,.5);
  \draw[<->,gainblue] (.5,.125)--(.5,.5);
  \node[anchor=west,align=left] at (.53,.32) {strict gap};
  \node[below] at (0,0) {$0$};
  \draw (.5,.012)--(.5,-.012) node[below] {$\tfrac12$};
  \draw (1,.012)--(1,-.012) node[below] {$1$};
  \node[anchor=north,align=center,text width=11cm] at (.54,-.16)
    {The linear entropy gain exceeds the quadratic action cost
     for a sufficiently small positive $t$.};
\end{tikzpicture}

\caption{The estimates used to create a strict remainder bound. The plot
uses $h_*=\mathcal E=1$ to illustrate their different orders in $t$.
For an actual packing, both constants are defined from its original
triangles. The difference between the upper line and lower curve is
$t h_*-t^2\mathcal E/2$.}\label{fig:weak-balance-gain}
\end{figure}

\section{The original filling and the actual hyperbolic volume}

The finite norm theorem of Appendix~\ref{app:uniform-filling-bound}
applies because $D_N$ has degree at most $d$ in each spectral variable.
For the original filling functionals it gives
\begin{equation}\label{wb:transfer}
 \left|\left(\bigotimes_e\mathcal L_N^{n_e}\right)D_N\right|
       \le\prod_e(1+|n_e|/2)\,\|D_N\|_{[-2,2]^R}.
\end{equation}
The twists are fixed, so this costs no exponential rate. The exact
finite identity also gives
$Z_N=(\bigotimes_e\mathcal L_N^{n_e})S_N$ and
$Z_{N,C}=(\bigotimes_e\mathcal L_N^{n_e})C_N$.

Appendix~\ref{app:morse-geometric-contribution} supplies, for
independently sufficiently long fixed signed powers,
\[
 Z_{N,C}=a_{N,\boldsymbol n}e^{N\mathcal A_{\boldsymbol n}},
 \qquad \frac1N\log|a_{N,\boldsymbol n}|\longrightarrow0,
\]
with nonzero amplitude for all sufficiently large levels, and
\begin{equation}\label{wb:filled-action}
 \Re\mathcal A_{\boldsymbol n}
   =\Phi-O_D\!\left(\sum_en_e^{-2}\right),\qquad
 2\pi\Re\mathcal A_{\boldsymbol n}
   =\Vol(S^3\setminus K_{\boldsymbol p}).
\end{equation}
The second identity uses the original primitive peripheral basis.
The actual periods satisfy the filling equation for
$\mu_e+p_e\lambda_e$, and the completed auxiliary end is a smooth
solid torus with cone angle $2\pi$. The main knot cusp stays complete.
The regulator's boundary correction cancels the filling correction
in the critical action. These are the geometric completion steps
proved there, not a new identification inferred from a shared
stationary equation.

Choose $0<\eta<\gamma_0$ and increase the common twist threshold
until the drop in~\eqref{wb:filled-action} is less than $\eta/2$.
Equations~\eqref{wb:remainder}--\eqref{wb:transfer} then give
\[
 \frac{Z_N}{Z_{N,C}}=1+O_{\boldsymbol n}(e^{-cN})
\]
for some $c>0$. The prescribed complete-color framing and component
factors have modulus one. In the full one-cut modified convention
the additional factor is $i^{N+3}/N$, which also has zero exponential
rate. Thus the same limit holds for the normalized invariant in
Theorem~\ref{wb:volume}, proving it. The argument also proves its
extension under~\eqref{wb:criterion}.

The theorem supplies an unchanged-diagram sufficient condition of
arbitrary finite size. It does not impose a tree or bridge-width
restriction. Whether a diagram satisfies the condition is a separate
geometric question; the theorem supplies neither a universal balance
claim nor a filling threshold uniform over growing diagrams.
